\documentclass[11pt]{article}
\title{Evaluating Agent Reliability Under Uncertain Outcomes and Ambiguous External Effects}
\date{}
\author{SIGMA AI-assisted research synthesis}
\usepackage[T1]{fontenc}
\usepackage[utf8]{inputenc}
\usepackage[margin=1in]{geometry}
\usepackage{amsmath,amssymb,array,booktabs,longtable,pdflscape}
\usepackage{microtype}
\usepackage{titlesec}
\usepackage{tikz}
\usetikzlibrary{fit}
\emergencystretch=8em
\raggedbottom
\linespread{1.07}
\usepackage{xcolor}
\definecolor{blue}{HTML}{24485E}
\titleformat{\section}{\Large\bfseries\sffamily\color{blue}}{\thesection}{0.8em}{}
\titleformat{\subsection}{\large\bfseries\sffamily}{\thesubsection}{0.7em}{}
\setlength{\parskip}{.6em}
\setlength{\parindent}{0pt}
\begin{document}
\maketitle
\begin{center}\fbox{\parbox{.92\linewidth}{\textbf{Needs review: TeX verification is incomplete.} Review the archived proofreading diagnostics before relying on the typesetting.}}\end{center}
\begin{center}\small\textit{AI-generated research synthesis \textperiodcentered{} SIGMA}\end{center}

\begin{abstract}
How should an agent\textquoteright{}s apparent completion be distinguished from a correct, recoverable, and reproducible outcome? This paper develops a bounded evaluation framework centered on fixed execution protocols, repeated task clusters, independently specified outcomes, and explicit observation horizons. Its strongest supported findings are analytical constraints on reliability claims rather than measured architectural improvements. Ambiguous acknowledgments can leave an agent unable to distinguish omission from an already committed external effect. However, issuing a retry does not imply committing a duplicate, and separating an effect from its completion marker can compromise safety or delivery in different ways.

These distinctions motivate separate assessments of external state, reporting correctness, service-contract conformance, and evaluator validity. Proposed statistical audits and recovery mechanisms remain conditional on unvalidated observation and implementation contracts. The resulting framework specifies what must be measured and controlled before successful demonstrations or recovery traces can support dependable-agent claims.
\end{abstract}
\section*{About this paper}
This paper was generated by SIGMA, an AI-assisted research system, from the project's recorded investigations, findings, and supporting evidence.
The analysis and conclusions reflect an automated synthesis of that research. Although SIGMA incorporates evidence tracking, critical review, and consistency checks, the resulting paper has not necessarily undergone independent verification or peer review.
Readers should evaluate its arguments, supporting evidence, and stated limitations before relying on its conclusions.
Research style: Distinguish algorithmic guarantees from tests, benchmarks and conjectures. Document counterexamples, complexity assumptions and reproducibility.
\section*{Evidence coverage limitation}
The evidence index omitted 32185 material saved research field fragment(s) across 36 source(s). These fragments were not reviewed for this paper. Interpretations that could depend on omitted material remain provisional; the complete frozen source snapshot and omission ledger are retained with the generation record.
\section{Introduction and problem statement}
Reliability requires an operational target rather than an informal completion label. The analysis defines success for a fixed task distribution, execution protocol, resource budget, rubric, and observation horizon; crashes and timeouts remain attempted executions. This prevents a claim about success per scheduled attempt from silently becoming a claim about accepted outputs or success after unrestricted retries.

External actions create a particularly consequential measurement problem: a local checkpoint or plausible acknowledgment need not reveal whether the intended business effect occurred. The focal task requires exactly one irreversible effect, making omission, duplication, truthful reporting, and safe escalation distinguishable outcomes rather than interchangeable notions of recovery.

The available evidence supports a focused methodological and analytical contribution, not a comprehensive account of agent reliability across domains. Benchmark-specific findings, broad failure-prevalence estimates, and causal improvements from architectural interventions remain unestablished; the article therefore concentrates on outcome measurement, ambiguous writes, and the contracts needed to evaluate recovery.

\section{Methods and evidence}
The synthesis compared the supplied definitions of execution-level reliability, stratified paired comparisons, and finite-roster auditing. It separated their target populations and sampling units, preserving the distinction between independent task clusters, potentially dependent runs within a cluster, and a roster fixed before audit selection.

The analytical review traced ambiguous-write histories through local intent, request delivery, durable effects, completion records, and acknowledgment handling. It retained the explicit dropped-retry counterexample and the correction distinguishing effect-first from marker-first split writes, rather than treating all retry or split-write behavior as equivalent.

The evidence assessment separated operational definitions, counterexamples, conditional invariant arguments, proposed experiments, and implementation evidence. No runtime results, independently validated service contracts, calibrated reference-error allowances, or measured architectural benefits were identified in the supplied selection.

\section{Reliability as a Fixed-Protocol Outcome}
Reliability is defined here as actual task success under a fixed execution protocol, rather than as an agent\textquoteright{}s declaration of completion. The protocol fixes the task distribution $D$, agent and model versions, prompts, tool interfaces, environment-generation policy, resource budget, success rubric, and observation horizon before outcomes are examined. These choices delimit the population of executions to which the reliability claim applies; changing versions, budgets, environments, or task distributions requires additional evidence rather than automatic transfer of the same claim. [C31]

For a task drawn from $D$, execute a prespecified batch of $r\geq 1$ scheduled runs and select an index $J$ uniformly from that batch. Let $Y_J\in\{0,1\}$ equal one exactly when the selected execution actually satisfies the independently specified end-state requirements and safety constraints by the deadline, and zero otherwise. The target is $p=\mathbb{E}[Y_J]=\mathbb{E}\!\left[\frac{1}{r}\sum_{j=1}^{r}Y_j\right]$. Uniform selection therefore defines reliability as the expected within-batch success fraction, without requiring the runs to be independent or to have identical outcome distributions. [C31]

Crashes and timeouts remain attempted executions and retain their scheduled places in the denominator. Deleting them would replace the intended target with success conditional on surviving or completing execution. Their outcome labels nevertheless follow the rubric: a timeout does not, by itself, establish whether the required external state was achieved, whereas failure to satisfy the rubric by its deadline makes the execution unsuccessful. Retaining an attempt and determining its actual outcome are thus separate requirements. [C31]

Per-attempt reliability also differs from accepted-output precision. If $Z_J$ denotes the automated acceptance decision, precision among accepted outputs concerns \$\textbackslash{}Pr(Y\_J=1\textbackslash{}mid Z\_J=1)\$, when that conditional probability is defined, rather than \$\textbackslash{}Pr(Y\_J=1)\$. Conditioning on acceptance excludes rejected attempts from the target population. The fixed-protocol definition instead asks whether an attempted execution meets the rubric, regardless of whether an automated checker accepts its result. [C31]

Likewise, a scheduled batch is not an instruction to keep trying until one execution succeeds. Eventual success after retries is a different target, and retry policies belong to the protocol whose reliability is being evaluated. Nor does uniform selection from a batch necessarily describe a fresh isolated execution: if the batch policy changes the selected run\textquoteright{}s outcome distribution, replacing the batch with a single run changes the experiment. A fresh-run interpretation is warranted only when the specified protocol supports it. [C31]

For comparisons, the target task population is represented by finitely many prespecified strata $h=1,\ldots,H$, with task distributions $D_h$ and weights $w_h\geq 0$ satisfying $\sum_{h=1}^{H}w_h=1$. In each positive-weight stratum, sample $n_h\geq 1$ task clusters; each cluster contains both systems\textquoteright{} scheduled batches on the same sampled task, with $r_h\geq 1$ runs per system. For system $s\in\{A,B\}$, define $p_{sh}$ as actual success probability for a task drawn from $D_h$ followed by uniform selection from its scheduled batch. The comparison target is $\Delta=\sum_{h=1}^{H}w_h(p_{Ah}-p_{Bh})$, so an advantage is defined relative to the same weighted task population. [C32]

The weighting operates first within a task and then across the target population. Writing $Y_{shij}$ for actual rubric satisfaction by system $s$ on run $j$ of task cluster $i$ in stratum $h$, the relevant task-level quantity is the batch average $\overline{Y}_{shi}=\frac{1}{r_h}\sum_{j=1}^{r_h}Y_{shij}$, whose expectation is $p_{sh}$. The fixed population weights, rather than the numbers of collected runs, determine each stratum\textquoteright{}s contribution to reliability. Pooling runs without preserving this hierarchy would generally produce a run-count-weighted mixture instead of the intended task-weighted target. [C32]

The design assumes independence between complete task clusters and identical sampling distributions within each stratum, while allowing unrestricted dependence among runs, systems, and judges inside a cluster. Paired systems are evaluated in isolated copies of the same initial environment so that one system does not alter the other\textquoteright{}s task state. These are requirements to establish in an evaluation, not properties supplied merely by calling a comparison paired; additional repetitions within a task do not create additional independent task samples.

The resulting standard for a reliability advantage is therefore a higher probability of actual rubric satisfaction across the specified weighted task population under the fixed protocol\textemdash{}not a larger collection of accepted outputs or repeated successes on the same tasks. [C32]

\textit{Claims discussed: }C31, C32
\section{What Statistical Evaluation Can Certify}
Statistical evaluation can certify a claim about actual success only when its sampling guarantees are connected to defensible outcome evidence. Automated acceptance and reference certification are observations about an execution, not interchangeable definitions of correctness. The central distinction is therefore between uncertainty from sampling executions and uncertainty about whether their labels represent successful outcomes; increasing the sample size addresses the former without necessarily resolving the latter. [C32, C33, C34]

For an execution selected under the fixed protocol, let $Y,O,R\in\{0,1\}$ denote actual rubric satisfaction, automated acceptance, and reference certification, respectively. Define \$p=\textbackslash{}Pr(Y=1)\$, \$q=\textbackslash{}Pr(O=1)\$, \$a=\textbackslash{}Pr(O=1,R=0)\$, and \$c=\textbackslash{}Pr(O=1,R=1,Y=0)\$. The reference-validity premise is $c\leq\eta$: an unconditional upper bound on jointly accepted failures in the target execution population, not the conditional error rate \$\textbackslash{}Pr(Y=0\textbackslash{}mid R=1)\$. A bound on \$\textbackslash{}Pr(R=1,Y=0)\$ would also control $c$, but imposes a stronger restriction than this argument requires. [C33]

Partitioning accepted failures by the reference judgment gives \$\textbackslash{}Pr(O=1,Y=0)\textbackslash{}leq a+\textbackslash{}eta\$. Since accepted true successes are a subset of all true successes, it follows that $p\geq\max\{0,q-a-\eta\}$, without assuming independence among truth, automated acceptance, and reference certification. False reference rejection can enlarge $a$ and make this lower bound conservative; false certification of actual failures instead threatens validity unless covered by $\eta$. When this localized allowance is the only constraint connecting judgments to truth, the sharp identified interval is $[\max\{0,q-a-\eta\},1]$: the observed judgments do not by themselves supply a nontrivial upper bound on actual success. [C33]

Judge agreement cannot establish the allowance needed for that conclusion. Consider two hypothetical worlds in which $O=R=1$ on every execution, but actual success is universal in one world and absent in the other. Their observed judgments are identical although $p=1$ in the first and $p=0$ in the second; additional judges sharing those outputs would not distinguish them. This is an identification counterexample, not a measured benchmark failure rate: excluding the unsuccessful world requires evidence connecting certification to actual outcomes. [C33]

A finite-sample certificate can then be constructed at the task-cluster level. For $N$ independent, identically distributed evaluation task clusters, each with $r$ scheduled runs, define $X_i=r^{-1}\sum_{j=1}^{r}O_{ij}$ and $\widehat q=N^{-1}\sum_{i=1}^{N}X_i$. Separately sample $m$ independent audit task clusters from the same task distribution, execute the same protocol, and select one run uniformly from each batch; averaging $Z_k=\mathbf{1}\{O_k=1,R_k=0\}$ gives $\widehat a$. The audit selection must not condition on automated acceptance for this estimator, and arbitrary dependence among repeated runs within a task remains permitted. [C32, C33]

Because the cluster observations are bounded in $[0,1]$, bounded-variable concentration and a union bound yield a conditional lower confidence certificate. For $0<\alpha<1$, write $t_N=\sqrt{\log(2/\alpha)/(2N)}$ and $t_m=\sqrt{\log(2/\alpha)/(2m)}$; then $L=\max\{0,\widehat q-t_N-\widehat a-t_m-\eta\}$ satisfies \$\textbackslash{}Pr(p\textbackslash{}geq L)\textbackslash{}geq1-\textbackslash{}alpha\$ under the stated sampling and reference-error assumptions. The probability statement concerns repeated sampling under that protocol, while $\eta$ remains a separate premise about the relationship between judgments and truth. In particular, the construction does not treat the $Nr$ executions as independent observations. [C32, C33]

An assumed sensitivity value for $\eta$ produces only a conditional sensitivity result. If an outcome-grounded procedure instead supplies an upper confidence bound $\eta_U$ for $c$ with failure probability at most $\beta$, substituting $\eta_U$ gives coverage of at least $1-\alpha-\beta$ by another union bound, without requiring independence between the procedures. No such calibrated allowance is supplied here; without a defensible restriction, $\eta=1$ is valid but forces $L=0$. Thus sampling precision cannot compensate for an unconstrained reference-validity error. [C33]

Comparative certification requires directional error control rather than two separate lower bounds on success. Let $S_s=O_sR_s$ denote joint acceptance and certification for system $s$; the exact decomposition is \$\textbackslash{}Pr(Y\_s=1)=\textbackslash{}Pr(S\_s=1)+\textbackslash{}Pr(Y\_s=1,S\_s=0)-\textbackslash{}Pr(Y\_s=0,S\_s=1)\$. For a lower bound on improvement of system $A$ over system $B$, the adverse errors are false certification of $A$ and missed true successes of $B$: the former can inflate the apparent intervention outcome, while the latter can depress the apparent baseline. These require unconditional, target-population allowances in their respective directions; controlling only false certification for both systems does not supply the needed upper control on baseline success. [C33]

The failure of lower-bound subtraction is visible in a simple counterexample. Suppose $S_A=1$ and $S_B=0$ on every paired execution, while $Y_A=Y_B=1$. The certification contrast is $1$, but actual improvement is $0$ because the baseline succeeds despite lacking certification. An allowance for missed true baseline successes addresses precisely this discrepancy; agreement among judges does not calibrate it. [C33]

For population comparison, the inferential unit is the complete paired task cluster underlying $\Delta=\sum_h w_h(p_{Ah}-p_{Bh})$. Clusters must be mutually independent and identically distributed within each prespecified stratum, while dependence among systems, scheduled runs, and judges inside a cluster is unrestricted. Averaging scheduled runs within a task preserves the intended task weighting and avoids claiming additional independent evidence merely from repetition. This design also requires distribution-matched sampling and noninterfering paired executions; shared state that couples different task clusters is not covered simply because within-cluster dependence is allowed. [C32]

Missing labels must weaken the comparison rather than silently change its target. For a lower improvement bound based on complete certification contrasts, assigning missing intervention certifications to $0$ and missing baseline certifications to $1$ decreases the observed contrast pointwise, even when missingness depends on outcomes. That conservative argument nevertheless presupposes well-defined complete outcomes and scheduled denominators: it does not justify deleting unexecuted attempts or inventing hypothetical results for them. Operational treatment of nonexecution must therefore be specified before outcomes are inspected. [C32, C33, C34]

Finite-roster auditing answers a different question from sampling independent task clusters, as distinguished in the table. Here every scheduled paired slot, its actual binary outcomes, and the observation cutoff are fixed before audit selection, and inspection must not alter those outcomes. Random audit selection, rather than an assumption of independent execution outcomes, supplies the sampling variation. A result about this roster does not establish reliability over a deployment task distribution without separate population-sampling and transport conditions. [C34]

Ambiguous evidence enters the finite-roster problem through sound sets of possible outcomes: missing information must widen the possibilities rather than be treated as favorable certification. Consequently, inspecting every roster entry can remove audit-selection uncertainty without necessarily resolving uncertainty about the outcomes themselves. This local treatment assumes that accepted evidence is sound; a globally invalid evidence premise cannot be repaired merely by marking individual records unknown.

The proposed enumeration and bound implementations remain untested, and the supplied derivations do not support a reproducible finite-roster coverage formula; dependence, masking, census cases, and confidence-level sensitivity remain implementation-validation requirements rather than reported results. [C34]

The statistical contribution is therefore a boundary on what successful evaluations can establish: confidence about observed acceptance becomes confidence about actual success only through explicit validity allowances and an appropriate sampling model. Comparative claims additionally need protection against missed baseline successes, while finite-roster claims remain tied to the audited executions. These distinctions replace an undifferentiated success score with a conditional certificate whose meaning depends jointly on the outcome evidence, the direction of judgment error, and the population to which the claim refers. [C32, C33, C34]

\begin{landscape}
\small
\begin{longtable}{@{}>{\raggedright\arraybackslash}p{\dimexpr(\linewidth-2\tabcolsep*4)/203*35\relax}>{\raggedright\arraybackslash}p{\dimexpr(\linewidth-2\tabcolsep*4)/203*42\relax}>{\raggedright\arraybackslash}p{\dimexpr(\linewidth-2\tabcolsep*4)/203*42\relax}>{\raggedright\arraybackslash}p{\dimexpr(\linewidth-2\tabcolsep*4)/203*42\relax}>{\raggedright\arraybackslash}p{\dimexpr(\linewidth-2\tabcolsep*4)/203*42\relax}@{}}
\caption{Two evaluation settings with different sources of sampling uncertainty. Their assumptions and inferential scope are not interchangeable.}\\
\toprule
\textbf{Setting} & \textbf{Source of sampling variation} & \textbf{Permitted outcome dependence} & \textbf{Essential validity conditions} & \textbf{Scope of the claim}\\
\midrule
\endfirsthead
\toprule
\textbf{Setting} & \textbf{Source of sampling variation} & \textbf{Permitted outcome dependence} & \textbf{Essential validity conditions} & \textbf{Scope of the claim}\\
\midrule
\endhead
\midrule
\multicolumn{5}{r}{\footnotesize Continued on next page}\\
\endfoot
\bottomrule
\endlastfoot
Independent task-cluster evaluation & Tasks and complete paired clusters sampled within prespecified strata & Unrestricted dependence within each cluster; independence required across complete clusters & Target-population matching, fixed weights and execution protocol, noninterference, and outcome-grounded directional judgment-error allowances & The prespecified task-population contrast under the fixed protocol; no automatic transport to changed environments, versions, or budgets\\
Fixed finite-roster audit & Audit selection from a complete roster whose outcomes are already fixed & Correlated execution outcomes are compatible with the setting; validity depends on the audit-selection design & Roster completeness, fixed cutoff and nonexecution rules, valid selection, nonreactive inspection, and sound possible-outcome sets & The scheduled finite roster; neither deployment-population reliability nor an empirically validated bound implementation follows\\
\end{longtable}
\par\smallskip{\footnotesize Evidence: [Research session 5, artifact 1], [Research session 5, artifact 3], [Research session 6, artifact 1], [Research session 7, artifact 1]; claims: C32, C33, C34}
\end{landscape}
\textit{Claims discussed: }C32, C33, C34
\section{Ambiguous Writes and the Limits of Local Recovery}
Ambiguous external writes expose an information limit on local recovery: the same pending checkpoint and timeout can require opposite actions to achieve exactly one external effect. The analytical example developed here locates that limit at the boundary between the agent\textquoteright{}s observations and the service\textquoteright{}s committed state. It does not, by itself, establish a reasoning deficit in the model or an empirical failure rate for deployed agents. [C35]

Consider a task requiring one irreversible external effect, such as appending one shipment order. Let $N$ count the effects attributable to that intended operation; correctness requires $N=1$, while omission and duplication are both incorrect. Ordinary requests have no built-in deduplication, and each committed request creates one additional effect. To isolate ambiguous commitment, restrict terminal recovery to skipping or issuing one retry, and initially assume that the retry, if chosen, certainly commits exactly one additional effect. [C35]

In history $W_0$, the initial request is lost before commitment, leaving $N=0$. In history $W_1$, the request commits once, leaving $N=1$, but its acknowledgment is lost. Both histories expose the same local intent record, timeout, messages, and other agent-accessible observations; no authoritative lookup or other distinguishing information is available before the recovery decision. A policy restricted to those inputs must therefore assign the same retry probability $q$ in both histories: random reasoning or replanning cannot distinguish histories whose accessible inputs and randomization distributions are identical. [C35]

The two histories nevertheless require opposite decisions. In $W_0$, skipping preserves the omission and retrying supplies the required effect, so the error probability is $e_0=1-q$. In $W_1$, skipping preserves the correct effect count and retrying creates a duplicate, so $e_1=q$. Consequently, $e_0+e_1=1$ and $\max(e_0,e_1)\geq 1/2$; choosing $q=1/2$ attains equality, whereas either deterministic choice fails with certainty in one history. This is a sharp worst-case bound within the restricted model, not a claim that the histories occur equally often in deployment. [C35]

The assured-commitment premise is essential to that calculation. In particular, the proposed extension \$\textbackslash{}Pr(\textbackslash{}text\{success\}\textbackslash{}mid W\_1)\textbackslash{}leq 1-q\$ is false under arbitrary relaxation of retry delivery. An always-retry policy whose retries all disappear has $q=1$, yet the original committed effect in $W_1$ remains the only effect, giving success probability one rather than the proposed upper bound of zero. This counterexample retracts the delivery relaxation, not the original conditional argument: choosing to retry and committing another effect are different events. [C36]

The bound also does not cover policies that obtain authoritative status information, escalate, compensate, wait indefinitely, or decline to claim completion. Those options change the available observations or terminal outcomes and cannot inherit the forced skip-or-retry result without a separate argument. Likewise, the worst-case calculation supplies neither deployment fault frequencies nor the relative harm of omission and duplication, so it does not select a practical retry policy. [C35]

To express these distinctions without conflating local progress with external completion, define the modeled state as $S=(E,D,L,Q,V,J)$. The accompanying state decomposition separates actual effects, service completion assertions, agent-visible state, network traffic, volatile execution, and evaluator-only history. For an intended operation with key $k$ and full payload $p$, correctness requires exactly one matching effect and no unintended effect; a completion entry in $D$ is an assertion to be checked against $E$, not an outcome oracle in its own right. [C37]

Operation identity must also remain distinct from delivery identity. A request envelope is $(x,k,p)$, where $x$ is fresh for each delivery attempt, while a retry preserves the operation key $k$ and payload $p$. An acknowledgment has its own fresh envelope identity, a correlation field referring to $x$, and a processing outcome. These definitions permit multiple attempts at one logical operation to be represented without treating either the number of requests or the number of receipts as the number of actual effects. [C37]

The transition semantics make loss and delay materially different. Sending adds a request to $Q$ without changing $E$ or $D$; dropping removes an envelope permanently, whereas delaying leaves it eligible for later delivery. Receipt destructively removes a request from the network pool into volatile service execution, with no retained broker lease or automatic redelivery; a service crash discards that volatile state but preserves $E$, $D$, and envelopes still in $Q$. Processing and acknowledgment enqueueing are separate events, so a committed effect can survive even when the corresponding reply never becomes available to the agent. [C37]

Restoring the agent changes $L$ to its saved state but does not reverse an irreversible effect or cancel outstanding network requests. Thus a restored pending operation can coexist with either no committed effect or an already committed effect, and a timeout alone cannot establish which external state holds. Treating resumed control flow as proof of completed external work can add a false completion claim to the underlying omission-or-duplication problem. [C35, C37]

The consequential implementation gap is observation equivalence: no executed harness supplied here demonstrates that the paired histories expose identical recovery inputs while an independent observer can distinguish their actual effects. Timing, tool metadata, restored memory, and outstanding requests must therefore be accounted for before the analytical pairing can support an experimental claim. Within its stated premises, however, the result changes the reliability argument decisively: successful local recovery is not sufficient evidence of the intended external outcome, and retry issuance must not be substituted for evidence of commitment. [C35\textendash{}C37]

\begin{longtable}{@{}>{\raggedright\arraybackslash}p{\dimexpr(\linewidth-2\tabcolsep*2)/104*20\relax}>{\raggedright\arraybackslash}p{\dimexpr(\linewidth-2\tabcolsep*2)/104*42\relax}>{\raggedright\arraybackslash}p{\dimexpr(\linewidth-2\tabcolsep*2)/104*42\relax}@{}}
\caption{State decomposition for the conditional ambiguous-write model. The components separate external outcomes from completion assertions and from information available during recovery.}\\
\toprule
\textbf{Component} & \textbf{Modeled contents} & \textbf{Recovery significance}\\
\midrule
\endfirsthead
\toprule
\textbf{Component} & \textbf{Modeled contents} & \textbf{Recovery significance}\\
\midrule
\endhead
\midrule
\multicolumn{3}{r}{\footnotesize Continued on next page}\\
\endfoot
\bottomrule
\endlastfoot
$E$ & Durable multiset of actual business effects, each containing an effect identifier, operation key $k$, and full payload $p$. & Defines the modeled external outcome; restoring agent state does not remove effects.\\
$D$ & Durable map from $k$ to a completed record containing the payload, result, and effect identifier. & Asserts completion but must be checked for correspondence with actual effects in $E$.\\
$L$ & Agent intent, saved state, and visible transcript. & May be restored locally without changing external durable state or outstanding network envelopes.\\
$Q$ & Separately accounted request and acknowledgment envelopes, partitioned into $Q_R$ and $Q_A$. & Survives service crashes, permits explicit drops, and is not a broker with automatic redelivery.\\
$V$ & Volatile service execution state, including in-flight workers. & Is discarded by a service crash; a request already consumed into this state is distinct from an envelope still in $Q$.\\
$J$ & Evaluator-only event history. & Supports reconstruction of modeled transitions without becoming agent-accessible recovery information.\\
\end{longtable}
\par\smallskip{\footnotesize Evidence: [Research session 4, artifact 2], [Research session 4, artifact 4]; claims: C37}
\textit{Claims discussed: }C35, C36, C37
\section{Recovery Contracts: Safety, Delivery, and Containment}
Recovery mechanisms support dependable-agent claims only when the property they protect is separated from the conditions under which they protect it. Preventing a second external effect, obtaining a first effect before a deadline, verifying that effect, and reporting completion truthfully are distinct obligations. The recovery contracts considered here therefore establish conditional constraints on execution, not interchangeable guarantees of task completion. Their comparison turns on operation identity, durable state, atomic enforcement, and the treatment of work that remains queued or resumes after recovery [C38\textendash{}C41].

Let $C(k,t)$ count the historical committed effects of logical operation $k$ by time $t$. At-most-once safety requires $C(k,t)\leq 1$ throughout execution, whereas effect completion by deadline $H$ requires $C(k,H)=1$; verified completion additionally requires appropriate authoritative evidence and a truthful report by that deadline. The clean-initialization arguments assume an initially unperformed operation, a stable identity, and a fixed payload preserved across the covered execution paths.

The supplied evidence consists of specifications, constructed schedules, and conditional deductions rather than executed implementations or measured reliability improvements; in particular, conformance of the actual effect boundary, durability, and universal execution checks remains unestablished [C40, C41].

Linked atomic effect-and-record creation addresses the interval in which an effect could otherwise become durable without its completion record. In the specified construction, serializable processing commits the covered effect and its linked record together, or rolls both back, and committed state survives the modeled crashes. A subsequent matching request can then return the existing result without creating another effect, provided the identity and suppressing state remain available. This protection concerns the modeled effect itself: atomic insertion of a dispatch intention does not establish atomicity with an external downstream outcome [C38\textendash{}C40].

Separating those writes creates different failure modes depending on their order. With effect-first processing, a crash after the effect commits but before its marker commits leaves an already performed operation without the state used to suppress another execution; a later effect-producing retry can therefore duplicate it. Under restricted marker-first suppression, a crash after the marker commits but before the effect commits can instead leave no effect while subsequent attempts are suppressed. The latter can preserve at-most-once safety while losing delivery, so splitting commits does not universally refute safety.

This distinction is specific to the restricted suppression behavior and does not establish the correctness of an additional repair path [C38].

Finite retention introduces a separate failure even when each effect-and-receipt commitment is atomic. In a constructed schedule with receipt lifetime $T=10$, the first effect commits at $t=1$, its acknowledgment is lost, and a recovered client sends a same-key retry at $t=2$ before stopping sends at $t=3$. The receipt expires at $t=11$, but the retry remains queued until processing at $t=12$, when the ordinary absent-receipt rule permits a second effect. Stable identity, atomic individual commitments, and a client send cutoff all hold in this schedule; what fails is coverage of delayed processing after expiry. Receipt absence can consequently no longer be equated with absence of a historical effect [C39].

A finite-retention safety argument must ensure that every later attempt either encounters the first receipt while it remains valid or is prohibited from creating an effect. One sufficient timing contract bounds the last effect-capable processing time, not merely request arrival: if all copies are issued by $S$ and each is permanently discarded or finally processed within $\Delta$ of issuance, then $L=S+\Delta$ bounds possible execution. Given a pre-established lower bound $a$ on first commitment, requiring $L<a+T$ keeps every possible later execution before the first receipt expires. This premise must cover queues, outages, redelivery, and resumed processing; average latency and a client retry limit do not establish it.

An alternative is a persistent tombstone checked atomically by every effect-producing path, although retaining suppression state does not itself retain the result needed for recovery [C39].

An immutable authorization interval provides a different conditional route to finite retention. Bind the operation and payload to $[a,D)$ with $0<D-a\leq T$, and require the service to reject effect creation at or after $D$, including work admitted earlier. If the first commitment occurs at $c\geq a$, its receipt remains valid until $c+T\geq D$: a duplicate before $D$ is blocked by the receipt, and one at or after $D$ is blocked by the authorization check. This argument requires nonrenewable authorization, clocks and retention behavior that survive restart and failover, and an eligibility check atomic with the actual effect; a paused worker cannot rely on authorization checked before the boundary.

It permits permanent omission when all requests arrive too late and supplies no completion-by-deadline guarantee [C39\textendash{}C41].

Cancellation instead requires a durable terminal decision about an operation. The proposed contract distinguishes PENDING, meaning no effect has committed, from APPLIED and CANCELLED, with effect creation and the APPLIED record committed indivisibly and successful cancellation acknowledged only after CANCELLED is durable. Its intended protection depends on ordering cancellation against execution at the same authoritative boundary, preserving terminal state through crashes, and binding responses to the correct operation. Stopping client issuance is not cancellation of queued requests, and a cancellation contract does not retroactively remove an effect that committed first.

Deleting cancellation state without a justified lifetime bound or another durable rejection rule reintroduces the retention obligation [C39, C40].

Epoch fencing targets stale execution under a different set of premises. The supplied safety argument combines durable linearizable epoch advancement, exact effect-count reads, and an append operation whose epoch check is atomic with effect creation. Its scope also requires at most one execution per invocation, serial fresh recovery rounds, and no resumption of abandoned controller flows; duplicate execution and overlapping controllers are therefore not silently covered. Durability of the epoch alone is insufficient if a worker can check an epoch and later create an effect without rechecking at the atomic boundary.

Under the stated service and controller premises, the argument supports an at-most-once invariant, not an unconditional recovery guarantee [C40, C41].

Progress requires additional premises beyond that invariant. The supplied progress discussion identifies limits on unsuccessful rounds and call durations as conditions for deterministic completion bounds, but fairness or average availability does not establish those limits. Similarly, a safe execution can produce no effect, or produce one effect without returning confirmation before the reporting deadline. These cases must remain distinct from verified completion, while an unresolved or indeterminate report must not be treated as a false success assertion merely because the task is incomplete [C41].

Incident containment addresses a different initial condition: excessive effects may already exist when recovery begins. The proposed durable closure mechanism prevents further covered effects only if closure is atomically ordered with the gate check and creation of every covered effect, persists through modeled failures and ordinary epoch changes, and is not reopened during the claimed interval. Because the historical effect count is monotone and is not reduced by refunds or corrective records, closure after duplication cannot restore an exactly-one-effect history. Nor does this conditional containment result establish bounded closure completion, notification delivery, or the safety of unmodeled downstream effects [C47].

The comparison therefore changes the central reliability claim from \textquotedblleft{}recovery succeeded\textquotedblright{} to a property-specific statement about what the service and controller jointly enforce. Linked atomic state, retained suppression, execution-time authorization, cancellation, fencing, and closure address different failure boundaries; none alone equates safe execution with delivered, verified, and truthfully reported completion. In particular, preserving safety by suppressing work can sacrifice delivery, while containing an incident can prevent additional effects without repairing historical correctness [C38\textendash{}C41, C47].

\begin{landscape}
\small
\begin{longtable}{@{}>{\raggedright\arraybackslash}p{\dimexpr(\linewidth-2\tabcolsep*4)/188*42\relax}>{\raggedright\arraybackslash}p{\dimexpr(\linewidth-2\tabcolsep*4)/188*42\relax}>{\raggedright\arraybackslash}p{\dimexpr(\linewidth-2\tabcolsep*4)/188*42\relax}>{\raggedright\arraybackslash}p{\dimexpr(\linewidth-2\tabcolsep*4)/188*42\relax}>{\raggedright\arraybackslash}p{\dimexpr(\linewidth-2\tabcolsep*4)/188*20\relax}@{}}
\caption{Conditional recovery contracts and their analytical boundaries. The entries compare protected properties and required premises, not measured performance or demonstrated implementation conformance.}\\
\toprule
\textbf{Mechanism} & \textbf{Conditional protection} & \textbf{Essential contract} & \textbf{Analytical failure or non-guarantee} & \textbf{Claim IDs}\\
\midrule
\endfirsthead
\toprule
\textbf{Mechanism} & \textbf{Conditional protection} & \textbf{Essential contract} & \textbf{Analytical failure or non-guarantee} & \textbf{Claim IDs}\\
\midrule
\endhead
\midrule
\multicolumn{5}{r}{\footnotesize Continued on next page}\\
\endfoot
\bottomrule
\endlastfoot
Linked atomic effect-and-record creation & Suppresses additional covered effects for the same operation while suppressing state remains available. & Stable identity and payload; serializable linked commitment of the actual covered effect and record; crash durability. & Atomicity over a dispatch intention does not cover an external downstream effect; no guarantee that a first effect or confirmation arrives. & C38, C39, C40, C41\\
Effect-first split writes & No general at-most-once protection under the specified retry behavior. & The effect commits before a separate suppressing marker. & A crash between writes leaves an existing effect unmarked; a later effect-producing retry can duplicate it. & C38\\
Restricted marker-first split writes & Can preserve at-most-once safety through suppression. & A durable marker precedes the effect and suppresses subsequent attempts under the restricted rule. & A crash between writes can leave zero effects with further delivery suppressed; an added repair path needs separate analysis. & C38\\
Finite receipts with a bounded execution horizon & Keeps all possible later processing within the first receipt's validity. & An all-copy last-processing bound $L$ and a first-commit lower bound $a$ satisfying $L<a+T$. & Send cutoffs, ingress bounds, and average latency do not bound queued or resumed execution. & C39\\
Persistent suppression tombstone & Blocks another covered effect after result-receipt expiry. & Tombstone installation atomic with the actual effect; durable persistence; mandatory checks on every covered path. & Suppression does not ensure result availability, first delivery, or successful reporting. & C39, C41\\
Immutable authorization interval & Blocks effect creation after $D$ while receipts cover earlier eligible execution. & Fixed interval $[a,D)$ with $0<D-a\leq T$; nonrenewable identity binding; atomic execution-time checks surviving failover. & Ingress-only checks or stale worker authorization are insufficient; late delivery can leave the operation permanently unperformed. & C39, C40, C41\\
Durable cancellation & Conditionally prevents covered execution after the cancellation terminal decision. & Authoritative ordering of cancellation and effect creation; durable terminal states; correct operation binding; retained rejection coverage. & Stopping sends is not cancellation; cancellation does not undo an effect ordered before it. & C39, C40\\
Epoch fencing & Supports at-most-once safety within the specified service and controller model. & Durable linearizable epoch advance; exact reads; atomic epoch check and effect creation; single execution per invocation; serial fresh rounds without abandoned-flow resumption. & Duplicate execution, concurrent controllers, and pre-existing corruption are outside the stated guarantee; progress needs additional timing and failure limits. & C40, C41\\
Durable closure & Prevents further covered effects after closure commits. & Closure atomically ordered with every covered effect path; persistence through failures and epoch changes; no reopening during containment. & Cannot reduce an already excessive historical effect count or establish correctness of compensation and unmodeled downstream consequences. & C47\\
\end{longtable}
\par\smallskip{\footnotesize Evidence: [Research session 7, artifact 5], [Research session 7, artifact 8], [Research session 6, artifact 4], [Research session 6, artifact 5], [Research session 6, artifact 6], [Research session 6, artifact 7], [Research session 7, artifact 2], [Research session 7, artifact 4], [Research session 7, artifact 7]; claims: C38, C39, C40, C41, C47}
\end{landscape}
\textit{Claims discussed: }C38, C39, C40, C41, C47
\section{Outcome Evidence and Evaluator Trust Boundaries}
Outcome certification requires an evidence contract separate from the recovery contract: the observer must establish what happened, what completion report was accepted, and how both relate to the evaluation cutoff without assuming that the recovery mechanism worked. In the state model, the durable effect multiset $E$ records business effects, whereas the completion map $D$ asserts completion; an entry in the latter is not an outcome oracle. The agent-visible state $L$ and evaluator-only journal $J$ therefore serve different evidentiary roles. The observer must inspect covered effects and audit their correspondence with completion records rather than substitute the records for the effects. [C37, C42]

The proposed observation boundary is deliberately narrow. For a single intended operation, the covered namespace begins empty, is isolated from other operations, and includes every immutable committed append effect covered by the rubric. Correct external state requires exactly one covered effect with the intended key and full payload, with no unintended covered effect. This boundary does not include unmodeled downstream consequences, and the ledger-backed specification supplies no general outcome oracle for open-ended agent work. Namespace isolation and complete effect capture are thus premises of the judgment, not consequences of finding one matching row. [C37, C42]

For an actual covered history $w$, let $E_H(w)=1$ mean that exactly one matching covered commit precedes cutoff $H$. Let $R_H(w)=1$ mean that exactly one final-report receipt precedes $H$ and that it is a well-formed, matching SUCCESS; malformed or additional final receipts violate this reporting contract. The horizon outcome is $Y_H(w)=E_H(w)\land R_H(w)$. A qualifying receipt denotes accepted reporting, not merely an agent's intention to report or evidence that it emitted a message. Consequently, a singleton-success judgment requires complete reporting capture or closure establishing content, uniqueness, and timeliness, as well as complete effect evidence. [C42]

The following schematic separates the agent's observations, service records, and proposed reference-observation boundary. Its capture arrows represent obligations of the observer design, not established independent capture paths; in particular, completion records remain objects to audit against effects. [C37, C42]
\$\textbackslash{}begin\{array\}\{rcl\}\textbackslash{}underbrace\{L,\textbackslash{} \textbackslash{}text\{received acknowledgments\}\}\_\{\textbackslash{}text\{agent-visible evidence\}\} \&\textbackslash{}longrightarrow\& \textbackslash{}text\{routine evidence \}d\textbackslash{} \textbackslash{}longrightarrow\textbackslash{} \textbackslash{}text\{mapper output \}m(d,a)\textbackslash{}\textbackslash{}\textbackslash{}underbrace\{E,\textbackslash{} D\}\_\{\textbackslash{}text\{durable service records\}\} \&\textbackslash{}xrightarrow\{\textbackslash{}text\{effect capture and correspondence audit\}\}\& \textbackslash{}text\{reference history \}w\textbackslash{}\textbackslash{}\textbackslash{}underbrace\{\textbackslash{}text\{accepted final receipts\},\textbackslash{} H,\textbackslash{} J\}\_\{\textbackslash{}text\{proposed observer boundary\}\} \&\textbackslash{}xrightarrow\{\textbackslash{}text\{receipt capture and ordering\}\}\& \textbackslash{}text\{reference history \}w\textbackslash{}end\{array\}\$

Ordering is substantive evidence rather than a timestamp-formatting detail. Commits, accepted receipts, and the cutoff require an enforceable common order; equal tick labels alone do not establish which event preceded another. Moreover, horizon correctness does not establish truthful completion at report-emission time: a commit can occur after the report but before the cutoff. A report-time judgment must inspect the commits preceding emission, while the horizon judgment inspects commits preceding the cutoff. An ordering ablation creates genuine uncertainty only if the retained evidence permits both relevant orders; removing an ordering certificate does not do so when retained, trusted timestamps already determine the order. [C42]

The mapper expresses evidence limitations through nonempty possible-value sets rather than forced binary labels. Given routine evidence $d$ and its current authorization or quarantine state $a$, it assigns component sets $U_E,U_R\subseteq\{0,1\}$ and returns $m(d,a)=\{e\land r:e\in U_E,\ r\in U_R\}$. A usable complete certificate determines its component by applying the full rubric to the certified contents. Without completeness, one matching effect or receipt cannot establish uniqueness; missing evidence ordinarily leaves $\{0,1\}$. Sound evidence of an irreversible violation, such as two distinct covered commits or an invalid final receipt, can nevertheless establish failure without complete history. [C42, C43]

Recognized contradictions require conservative widening rather than an empty outcome set. Incompatible exact-count certificates for the same namespace and cutoff cannot both be true, but that contradiction does not itself identify the actual count. Every affected component therefore becomes $\{0,1\}$; a contradiction in shared cutoff ordering can affect both components. If independently sound evidence still establishes failure of one component, the conjunction can remain $\{0\}$ despite uncertainty in the other. Missing or contradictory evidence is thus neither successful completion nor evidence that no external effect occurred. [C42, C43]

Three judgments must remain separate: whether the mapper follows its specification for the supplied evidence version, whether the evidence's factual assertions are valid, and whether the returned set contains the actual outcome. The containment argument is conditional. If every factual constraint used by the mapper is true, complete where claimed, correctly scoped, and correctly ordered, and the mapper conforms, then each component set contains its actual value; applying conjunction yields $Y_H(w)\in m(d,a)$. Conservative inclusion of impossible component combinations only enlarges the set.

Conformance alone, however, does not establish this result: a false completeness assertion can make a correctly implemented mapper return $\{1\}$ for an actual failure. [C42, C43]

The analytical fixtures in the table expose this attribution distinction without claiming executed tests. A sound held-out reference can reveal that a routine certificate omitted a second effect or an additional final receipt, but that disclosure does not retroactively change whether the mapper followed its specification on the original input. Conversely, actual-outcome containment can occur despite invalid evidence, as when conservative widening includes both binary outcomes. Mapper conformance, evidence validity, and containment therefore require separate recorded judgments rather than a universal rejection rule based solely on exclusion of the actual outcome. [C43]

Certificate withdrawal adds an authorization judgment, not another task-outcome label. An unresolved contradiction or disclosed breach requires a hold on certification dependent on the affected trust domain; the hold is neither proof of task failure nor a repair of false evidence. The original evidence and output should remain distinguishable from a new evidence version produced after quarantine or replacement. Widening the new version can restore conservative containment, but it does not establish the soundness of the original issuer or of certificates outside the immediately affected case. Defensible withdrawal scope and independently justified hold-clearing criteria remain necessary. [C42, C43]

The consequential applicability gap is the validity of the reference observer itself. Complete durable-effect capture, accepted-receipt capture, cutoff serialization, isolation, crash behavior, bypass prevention, and nonexecution evidence have not been established by an inspected implementation or executed pilot. Separate retention paths and archive digests do not establish these properties: routine and reference outputs can agree while sharing a storage, extraction, or reporting failure that omits the same event. The proposed guarantees therefore depend on validating the underlying observation contracts, not merely comparing the resulting labels. [C42]

This separation changes the reliability claim that an evaluation can support. A recovery trace, a completion record, and a conforming outcome mapper can coexist with an incorrect external outcome or unsound certification. Evidence for dependable completion must consequently connect covered external state and qualifying reporting to a valid observation boundary, while attributing mapping defects, evidence-contract breaches, and task failures separately. The decisive inference is not that the system produced a completion claim, but that the evidence supporting that claim is adequate to contain the actual outcome. [C37, C42, C43]

\begin{quote}
\textbf{Diagram (readable relationship list).} Proposed separation of agent-visible completion evidence from actual effects and reference adjudication. Completion records and reports are inputs to evaluation, not substitutes for an independently validated account of the covered external outcome. The boundaries indicate intended responsibilities, not demonstrated independence or implementation conformance.
\begin{description}
\item[\textbf{Node 1.}] Agent intent, checkpoint, and transcript
\item[\textbf{Node 2.}] Effect-producing service
\item[\textbf{Node 3.}] Durable actual effects $E$
\item[\textbf{Node 4.}] Durable completion records $D$
\item[\textbf{Node 5.}] Terminal report and receipt evidence
\item[\textbf{Node 6.}] Routine evidence mapper
\item[\textbf{Node 7.}] Proposed held-out reference observer
\item[\textbf{Node 8.}] Separate outcome, mapping, and validity judgments
\item[\textbf{Edge 1.}] Node 1 points to node 2: Requests and retries with stable operation identity
\item[\textbf{Edge 2.}] Node 2 points to node 3: Creates covered business effects
\item[\textbf{Edge 3.}] Node 2 points to node 4: Maintains completion assertions
\item[\textbf{Edge 4.}] Node 4 points to node 1: Acknowledgment or read-back evidence, when available
\item[\textbf{Edge 5.}] Node 1 points to node 5: Submits a terminal report
\item[\textbf{Edge 6.}] Node 4 points to node 6: Routine service evidence
\item[\textbf{Edge 7.}] Node 5 points to node 6: Routine reporting evidence
\item[\textbf{Edge 8.}] Node 3 points to node 7: Requires validated complete effect capture
\item[\textbf{Edge 9.}] Node 5 points to node 7: Requires validated receipt and cutoff capture
\item[\textbf{Edge 10.}] Node 6 points to node 8: Mapped labels and unresolved evidence
\item[\textbf{Edge 11.}] Node 7 points to node 8: Reference evidence with explicit validity obligations
\item[\textbf{Boundary 1.}] Agent-controlled execution state: node 1
\item[\textbf{Boundary 2.}] Service contract requiring implementation validation: node 2, node 3, node 4
\item[\textbf{Boundary 3.}] Proposed evaluation and adjudication: node 5, node 6, node 7, node 8
\end{description}
\par\smallskip{\footnotesize Evidence: [Research session 4, artifact 2], [Research session 7, artifact 6], [Research session 7, artifact 8], [Research session 7, artifact 9]; claims: C37, C42, C43}
\end{quote}
\begin{landscape}
\small
\begin{longtable}{@{}>{\raggedright\arraybackslash}p{\dimexpr(\linewidth-2\tabcolsep*4)/141*42\relax}>{\raggedright\arraybackslash}p{\dimexpr(\linewidth-2\tabcolsep*4)/141*20\relax}>{\raggedright\arraybackslash}p{\dimexpr(\linewidth-2\tabcolsep*4)/141*20\relax}>{\raggedright\arraybackslash}p{\dimexpr(\linewidth-2\tabcolsep*4)/141*35\relax}>{\raggedright\arraybackslash}p{\dimexpr(\linewidth-2\tabcolsep*4)/141*24\relax}@{}}
\caption{Conditional analytical fixtures separating mapper conformance, evidence validity, and actual-outcome containment. Assume an initially authorized issuer and a mapper that implements the specified rules. Clean routine certificates assert exactly one matching effect and exactly one qualifying final receipt before the cutoff. These are reconstructed fixtures, not executed observations or reproductions of unavailable historical traces.}\\
\toprule
\textbf{Evidence condition and actual history} & \textbf{Routine outcome set} & \textbf{Mapper conforms} & \textbf{Evidence assertions valid} & \textbf{Actual outcome contained}\\
\midrule
\endfirsthead
\toprule
\textbf{Evidence condition and actual history} & \textbf{Routine outcome set} & \textbf{Mapper conforms} & \textbf{Evidence assertions valid} & \textbf{Actual outcome contained}\\
\midrule
\endhead
\midrule
\multicolumn{5}{r}{\footnotesize Continued on next page}\\
\endfoot
\bottomrule
\endlastfoot
Complete, correct certificates; one matching effect and exactly one qualifying final receipt precede the cutoff. & $\{1\}$ & Yes & Yes & Yes\\
Recognized certificates make incompatible exact-count claims for the same namespace and cutoff; the affected effect component is widened. & $\{0,1\}$ & Yes & No: at least one assertion is false & Yes\\
Apparently complete clean certificates omit a second covered commit before the cutoff; the actual task outcome is failure. & $\{1\}$ & Yes & No & No\\
Apparently complete clean certificates omit an additional final receipt before the cutoff; the actual task outcome is failure. & $\{1\}$ & Yes & No & No\\
\end{longtable}
\par\smallskip{\footnotesize Evidence: [Research session 7, artifact 9], [Research session 7, artifact 6]; claims: C42, C43}
\end{landscape}
\textit{Claims discussed: }C37, C42, C43
\section{Discriminating Experiments and the Supported Reliability Claim}
The smallest coherent validation program separates three questions: whether the service implements the recovery contract, whether the observer can establish the relevant outcomes, and whether a recovery policy improves agent performance under controlled conditions. These questions require distinct tests because a favorable policy comparison cannot by itself validate either the underlying service or the evaluator. The supplied evidence establishes analytical constraints and experimental obligations, but not an executed fault-injection study, a validated recovery implementation, calibrated judge-error allowances, or a measured architectural benefit.

The contribution is therefore a discriminating validation program, not an implementation result. [C44, C45]

Service-contract testing should begin with deterministic replay of the specified transition traces, checking the durable effect multiset and completion-record map after every event and accounting for every pending request and acknowledgment. Crash injection should distinguish failure before request receipt, after receipt but before processing, after processing, and after acknowledgment enqueueing; these boundaries test different loss semantics rather than interchangeable instances of a generic timeout. The audit must establish whether the business effect actually lies inside the claimed atomic boundary, alongside durability, serialization, retention, collision prevention, and payload-mismatch rejection.

Neither an API name nor the existence of a key table establishes those properties. [C44, C45]

The harness also needs tests of its intended information constraints. Paired histories designed to expose the same recovery information must be checked for equality of agent-visible observations, including possible timing, tool-metadata, restored-memory, and shared-state differences. For fencing-based recovery, the service tests should additionally exercise lost and delayed epoch-advance replies, lost append replies, lost confirmations, and restart behavior that might reuse epochs or resume abandoned control flows.

Duplicate delivery and overlapping controllers should be labeled separately as stress conditions outside the stated single-execution and serial-recovery premises, rather than interpreted as counterexamples to guarantees that exclude them. [C44, C45, C46]

Observer-validity testing should proceed independently of those service tests. An actual pilot must compare captured service histories and report receipts against independently retained outcomes, exercising cutoff ordering, namespace isolation, crash behavior, bypass prevention, and launcher evidence for nonexecution. Blinded cases should preserve routine evidence separately from reference evidence, with outputs recorded before reference disclosure. Missing, contradictory, and silently false evidence are discriminating conditions: they test whether a failure arises in the mapping implementation, in an evidence contract, or in end-to-end outcome certification.

A mapper that follows its specification should not be classified as incorrectly implemented merely because its inputs falsely describe the world. [C44, C45]

The table identifies the different conclusions these tests can support. Finite agreement on synthetic or replayed cases supports only the checked implementation behavior; it does not establish universal service conformance or real-world observer soundness. Statistical audit implementations likewise require checks under perfect within-task dependence, asymmetric judge errors, and outcome-dependent evidence loss, as well as boundary cases such as all-unknown evidence, census audits, and rejection exhaustion. These are implementation-validation obligations, distinct from the assumptions of any analytical coverage result. [C44, C45]

Only after service and observation checks should the program compare recovery policies in repeated, task-blocked trials. The core comparison should include fencing, read-back-only retry, restart, and safe escalation; continued execution provides a further comparator for assessing recovery against proceeding without restart. Each policy should receive isolated copies of the same initial environment, with matched model and prompt versions, tool access, execution budgets, and environment snapshots. Condition order should be randomized, namespaces isolated, and actual effects and rejection records independently audited.

These controls are intended to distinguish recovery-policy behavior from differences in capability, resources, or external state; no effect size is supplied by the present evidence. [C44, C45]

Reproducibility also requires measuring whether the agent preserves the intended operation key and full payload across retries, restoration, replanning, and concurrent work. The task distribution, stratum weights, scheduled repetitions, success rubric, observation horizon, and comparison should be fixed before inspecting outcomes. Repeated runs belong within task clusters, with uncertainty estimates respecting within-task dependence rather than treating every run as an independent observation. Nonexecution and budget exhaustion must remain in the scheduled-run denominator, since excluding them would answer a different question from whether the policy completes its assigned executions. [C44]

The outcome report should retain a vector of results rather than collapse recovery into a single success label. At minimum, it should distinguish duplicate effects, omitted effects, actual correctness, verified completion, false completion reports, unresolved operations, safe escalation or termination, and recovery outcomes, alongside latency, service calls, and inclusive cost. Severity and additional effects after reported completion require separate measurement. In particular, escalation or containment must not be counted as successful task completion merely because further harm was avoided.

Deployment fault frequencies and asymmetric omission and duplication costs would also need estimation before these comparisons could justify a preferred policy. [C41, C44]

Deadline claims require an additional evidential step beyond observing safe executions. At-most-once safety, effect completion by a deadline, verified completion, and truthful reporting remain distinct properties: preventing another effect does not establish that the required effect occurs or that confirmation reaches the client in time. The saved progress argument requires justified limits on unsuccessful rounds and call duration, which fairness or average availability alone does not supply. Fault-injection tests should therefore compare observed calls and durations with the proposed bounds, while keeping empirical test results separate from deterministic guarantees whose environmental premises still require justification. [C41, C44, C45]

The first supported scope remains a single ledger-backed operation with exactly one required effect, not open-ended agent work or dependent multi-service execution. Held-out tasks, changed versions, budgets, and environments require new evidence or separately justified transport conditions; persistent shared state and dependent workflows require additional treatment. Delayed harm, persistent corruption, irreversible downstream consequences, and recovery quality also exceed a fixed-horizon binary outcome. In particular, a finite observation horizon cannot establish the absence of all future late effects without a verified fencing contract.

No current benchmark-specific conclusion follows from this program, and unavailable original arguments and omitted obligations prevent claims that the proposed fixtures exhaust or reproduce all earlier cases. [C46]

The evidence thus changes the reliability argument from an architectural preference into a set of distinguishable claims with different validation requirements. Service conformance, observer validity, and controlled policy benefit must each be established on their own terms; none can be inferred from a successful recovery trace or agreement among evaluators. What is supported is a bounded analytical framework that identifies what a reliability improvement would have to demonstrate\textemdash{}not a finding that any proposed architecture already delivers that improvement. [C41, C44, C45, C46]

\begin{landscape}
\small
\begin{longtable}{@{}>{\raggedright\arraybackslash}p{\dimexpr(\linewidth-2\tabcolsep*3)/165*39\relax}>{\raggedright\arraybackslash}p{\dimexpr(\linewidth-2\tabcolsep*3)/165*42\relax}>{\raggedright\arraybackslash}p{\dimexpr(\linewidth-2\tabcolsep*3)/165*42\relax}>{\raggedright\arraybackslash}p{\dimexpr(\linewidth-2\tabcolsep*3)/165*42\relax}@{}}
\caption{Discriminating validation layers and the conclusions their results would permit. These are proposed tests, not completed experiments.}\\
\toprule
\textbf{Validation layer} & \textbf{Discriminating evidence} & \textbf{Interpretation of a negative result} & \textbf{Exact evidence source}\\
\midrule
\endfirsthead
\toprule
\textbf{Validation layer} & \textbf{Discriminating evidence} & \textbf{Interpretation of a negative result} & \textbf{Exact evidence source}\\
\midrule
\endhead
\midrule
\multicolumn{4}{r}{\footnotesize Continued on next page}\\
\endfoot
\bottomrule
\endlastfoot
Service and recovery contract & Per-event effect and completion-record checks; complete request accounting; crash-boundary replay; independent inspection of actual atomicity and durability. & A violated service premise limits application of the conditional recovery guarantee; a scope-violating stress case must be identified separately. & [Research session 4, artifact 4]; [Research session 7, artifact 2]\\
Observer and evidence mapping & Blinded routine outputs compared with independently retained effects and receipts; missing, conflicting, and falsely complete evidence cases. & Distinguish mapper implementation failure from evidence-contract failure and end-to-end certification failure; a shared observation failure cannot be repaired merely by relabeling local evidence as unknown. & [Research session 7, artifact 9]\\
Task-blocked recovery-policy comparison & Matched tasks, versions, snapshots, access, and budgets; randomized condition order; isolated namespaces; independent effect audits; retained nonexecution and exhaustion. & Report the observed safety, completion, recovery, latency, and cost contrasts without converting them into an unsupported architectural or deployment-wide benefit. & [Research session 7, artifact 2]; [Research session 5, artifact 3]\\
\end{longtable}
\par\smallskip{\footnotesize Evidence: [Research session 4, artifact 4], [Research session 7, artifact 2], [Research session 7, artifact 9], [Research session 5, artifact 3]; claims: C44, C45, C46}
\end{landscape}
\textit{Claims discussed: }C41, C44, C45, C46
\section{Algorithms, results, and limitations}
\subsection{C31 -- Established}
The supplied analysis adopts a fixed-protocol reliability target: draw a task from $D$, execute its prespecified batch of $r\geq 1$ runs, select a run $J$ uniformly, and set $p=\mathbb{E}[Y_J]$, where $Y_J$ records actual rubric satisfaction by the deadline. Crashes and timeouts remain attempted executions.
This is an explicit operational definition, not an empirical estimate. Fixing versions, tools, budgets, environment policy, and the observation horizon bounds the population to which a later estimate would apply. Uniform selection within a scheduled batch does not automatically describe a fresh isolated execution or success after unlimited retries.
\textit{Evidence: }[Research session 1, artifact 1], [Research session 4, artifact 1]
\subsection{C32 -- Established}
The paired comparison design targets $\Delta=\sum_h w_h(p_{Ah}-p_{Bh})$ over prespecified task strata. It assumes independent task clusters while allowing unrestricted dependence among runs, systems, and judges inside each cluster.
The saved definitions make the task cluster the independent sampling unit and average scheduled runs within that unit. This preserves the intended task weights rather than allowing tasks with more runs to dominate the target. It specifies a design; it does not establish that an implemented evaluation satisfies independence or noninterference.
\textit{Evidence: }[Research session 5, artifact 1], [Research session 5, artifact 3]
\subsection{C33 -- Unresolved}
The proposed connection between automated acceptance, reference certification, and actual success requires explicit reference-validity assumptions. In the single-system formulation, the allowance \$c=\textbackslash{}Pr(O=1,R=1,Y=0)\textbackslash{}leq\textbackslash{}eta\$ is an unconditional target-population probability, not the conditional error rate \$\textbackslash{}Pr(Y=0\textbackslash{}mid R=1)\$.
The quantities are explicitly distinguished in the supplied formulation, but a defensible allowance is not supplied. Comparative claims additionally require directional truth-error control, including missed true baseline successes. Agreement among judges does not calibrate these allowances.
\textit{Evidence: }[Research session 4, artifact 1], [Research session 5, artifact 1], [Research session 5, artifact 3]
\subsection{C34 -- Unresolved}
Finite-roster auditing is specified as a separate evaluation setting in which scheduled paired slots and their binary outcomes are fixed before audit selection. Its application depends on roster completeness, valid selection, and sound possible-outcome sets; the proposed enumeration and bound implementations remain untested.
The fixed-roster assumptions differ from sampling independent task clusters from a target distribution. The excerpts identify masking, missing evidence, census audits, correlated outcomes, and randomization as implementation obligations. They do not expose enough of the complete statistical argument to reproduce or certify a coverage formula.
\textit{Evidence: }[Research session 6, artifact 1], [Research session 7, artifact 1], [Research session 7, artifact 9]
\subsection{C35 -- Established}
The ambiguous-write model contains two histories with the same agent-accessible recovery information but different external outcomes: loss before commitment and commitment followed by acknowledgment loss. Under the restricted choices of skipping or making one retry that certainly commits, the two histories require opposite recovery choices to attain exactly one effect.
If the initial request did not commit, skipping leaves zero effects and the stipulated committing retry supplies the required effect. If it already committed, skipping preserves one effect while the stipulated committing retry creates another. This is a conditional analytical incompatibility; identical observations in an actual harness have not been demonstrated.
\textit{Evidence: }[Research session 1, artifact 2], [Research session 4, artifact 3], [Research session 6, artifact 2]
\subsection{C36 -- Established}
The inequality \$\textbackslash{}Pr(\textbackslash{}text\{success\}\textbackslash{}mid W\_1)\textbackslash{}leq 1-q\$, where $q$ is retry probability in an already-committed history $W_1$, does not survive arbitrary relaxation of the retry-commitment premise. An always-retry policy whose retries all disappear has $q=1$ while preserving exactly one committed effect.
The counterexample separates a policy decision from network delivery and service commitment. Its already-committed effect remains present even though no retry takes effect, contradicting the proposed zero upper bound. It retracts the delivery relaxation, not the original argument under assured retry commitment.
\textit{Evidence: }[Research session 6, artifact 1], [Research session 6, artifact 2]
\subsection{C37 -- Established}
The proposed recovery state machine distinguishes the durable actual-effect multiset $E$, durable completion map $D$, agent state $L$, network envelopes $Q$, volatile service state $V$, and evaluator-only history $J$. A completion entry in $D$ asserts completion but is not itself the outcome oracle.
These are explicit model definitions. Requests have fresh delivery identities while retries preserve the operation key and payload; request and acknowledgment envelopes are accounted for separately. Distinguishing these objects makes it possible to specify crash and loss schedules without silently equating a receipt, a checkpoint, and a business effect.
\textit{Evidence: }[Research session 4, artifact 2], [Research session 4, artifact 4]
\subsection{C38 -- Established}
The saved correction distinguishes split-write failure modes: effect-first splitting permits duplication, whereas restricted marker-first suppression can preserve at-most-once safety while losing delivery. Therefore, splitting an effect and its marker does not universally refute at-most-once safety.
The correction explicitly withdraws an unqualified safety claim and separates the property threatened by each ordering. The article can state that distinction and identify the relevant crash interval, but the complete saved schedules must be inspected before presenting an exhaustive transition-level proof.
\textit{Evidence: }[Research session 7, artifact 5], [Research session 7, artifact 8]
\subsection{C39 -- Unresolved}
A recovery argument relying on indefinitely available completion records does not transfer unchanged to finite retention. The proposed alternatives require justified execution horizons, persistent tombstones, or service-enforced rejection or fencing that covers queued and resumed work.
The source excerpts explicitly identify record removal as outside the earlier invariant's assumptions. Average latency and client-side retry limits do not establish a last-processing bound for delayed requests. Execution-time authorization checks and retention behavior remain implementation obligations rather than demonstrated safeguards.
\textit{Evidence: }[Research session 6, artifact 4], [Research session 6, artifact 5], [Research session 6, artifact 7]
\subsection{C40 -- Unresolved}
Cancellation, immutable authorization intervals, and epoch fencing are conditional recovery candidates, not interchangeable proven repairs. Their stated guarantees depend on stable identity, durable state, authoritative reads where required, and atomic enforcement at every covered effect-producing path.
The alternatives impose different concurrency, restart, retention, and timing assumptions. In particular, the epoch-based argument is scoped to serial fresh recovery rounds and suppression of abandoned controller flows. No supplied implementation establishes conformance to the complete contracts.
\textit{Evidence: }[Research session 6, artifact 6], [Research session 6, artifact 7], [Research session 7, artifact 2], [Research session 7, artifact 4]
\subsection{C41 -- Established}
At-most-once safety, effect completion by a deadline, verified completion, and truthful reporting are distinct properties. The saved progress discussion additionally requires justified unsuccessful-round and call-duration limits; fairness or average availability alone does not establish those deterministic premises.
The outcome definitions distinguish zero effects, one effect, duplication, safe escalation, and evidence-backed reports. A bound preventing additional effects cannot alone show that an effect occurs or that confirmation reaches the client by the deadline. The excerpts identify, but do not empirically justify, the additional progress assumptions.
\textit{Evidence: }[Research session 5, artifact 2], [Research session 6, artifact 7], [Research session 7, artifact 4]
\subsection{C42 -- Unresolved}
The proposed outcome observer requires complete covered-effect capture, accepted report-receipt capture, cutoff ordering, namespace isolation, and explicit treatment of missing or contradictory evidence. Sound certificates and sound possible-outcome sets remain assumptions requiring validation.
The observer must establish the task outcome without assuming the treatment's intended invariant. Unknown evidence cannot silently become successful completion or evidence of no effect. Shared storage, extraction, or reporting failures can undermine apparent independence even when routine and reference outputs agree.
\textit{Evidence: }[Research session 7, artifact 1], [Research session 7, artifact 6], [Research session 7, artifact 8], [Research session 7, artifact 9]
\subsection{C43 -- Established}
Mapper conformance and outcome validity require separate judgments. A mapper can follow its specification even when unsound or incomplete supplied evidence excludes the actual outcome; that observation challenges evidence validity rather than automatically demonstrating a mapper implementation failure.
The supplied correction identifies this attribution error explicitly. Separate labels are needed for mapping behavior, reference validity, service or controller contract failures, and exposed trust violations. The bounded extract does not contain the complete historical cases needed to reproduce every classification.
\textit{Evidence: }[Research session 7, artifact 9]
\subsection{C44 -- Unresolved}
The proposed discriminating evaluation compares recovery policies under controlled schedules and matched tasks, versions, tool access, environment snapshots, and budgets. It calls for isolated namespaces, randomized condition order, independent effect audits, inclusion of nonexecution and budget exhaustion, and separate safety, completion, recovery, latency, and cost outcomes.
These controls are saved experimental obligations, not completed experiments. Comparing fencing, read-back-only retry, restart, continued execution, and safe escalation would distinguish protocol behavior from differences in model capability or resources. No effect size or reliability-cost frontier can yet be reported.
\textit{Evidence: }[Research session 4, artifact 2], [Research session 5, artifact 3], [Research session 7, artifact 2]
\subsection{C45 -- Limitation}
The supplied evidence selection does not establish an executed fault-injection study, a validated recovery implementation, calibrated judge-error allowances, or a measured architectural reliability benefit. Recorded arguments and agreement are not independent verification.
The evidence repeatedly leaves execution, independent auditing, service conformance, and error calibration open. Several records explicitly state that no solution or complete verification is claimed. The article must therefore report methodological results and scoped counterexamples without converting planned tests into observations.
\textit{Evidence: }[Research session 1], [Research session 4], [Research session 5], [Research session 6], [Research session 7], [Research session 5, artifact 3]
\subsection{C46 -- Limitation}
The supported scope is narrower than the project's full reliability agenda: exactly-one-effect ledger-backed operations do not establish correctness for open-ended tasks, multi-service workflows, delayed harms, or current benchmarks. Several source extracts also omit original arguments and historical obligations.
The source notes explicitly exclude broader outcome-oracle and downstream-effect questions and retain benchmark auditing as future work. Missing originals and omitted obligations prevent certification of complete historical compatibility or coverage. Operational preflight and context failures are not benchmark measurements of deployed agents.
\textit{Evidence: }[Research session 4, artifact 2], [Research session 4, artifact 4], [Research session 5, artifact 2], [Research session 7, artifact 1], [Research session 7, artifact 9]
\subsection{C47 -- Unresolved}
Incident containment is distinct from restoring historical correctness. The proposed durable closure mechanism addresses further covered effects only if closure is atomically ordered with every covered effect path; an already excessive monotone effect count lies outside clean-initialization guarantees.
The incident formulation counts historical primary effects without subtracting refunds or corrective records. Preventing further effects therefore cannot make an already duplicated history satisfy an exactly-one-effect rubric. Closure authorization, persistence, bypass prevention, and downstream coverage remain implementation obligations.
\textit{Evidence: }[Research session 7, artifact 7]
\section{Limitations}
\begin{itemize}
\item The contribution is a bounded analytical and evaluation plan. The supplied selection establishes neither implementation conformance nor empirical superiority of a recovery architecture.
\item The mathematical material includes definitions, local counterexamples, and conditional proof candidates. Complete coverage calculations and several transition-level arguments are not available in the preview and must not be reconstructed as validated theorems.
\item A ledger-backed exactly-one-effect task provides a narrow outcome model. It does not resolve open-ended correctness, payload adequacy beyond the rubric, downstream delayed harm, or multi-service task correctness.
\item Reference validity remains a material assumption. Agreement among automated judges, critics, or synthetic observers cannot establish independent ground truth or calibrate directional truth-error allowances.
\item The bounded selection omits original arguments and historical obligations. It cannot support an assertion of exhaustive source review, complete reconciliation across branches, or resolution of the omitted duties.
\item No benchmark-specific measurements, literature comparison, deployment-weighted failure rates, latency results, cost measurements, or controlled architectural effect sizes are supplied.
\end{itemize}
\section{Open obligations}
\begin{itemize}
\item Recover the complete statistical and recovery arguments before asserting exact coverage formulas, exhaustive schedule results, or compatibility across the proposed evaluation settings.
\item Calibrate target-population reference-error allowances using independent outcome evidence, including baseline false rejection in paired comparisons; otherwise retain those allowances as sensitivity parameters.
\item Implement and audit task-cluster and finite-roster calculations separately. Exercise dependence, shared errors, masking, outcome-dependent evidence loss, all-unknown evidence, census audits, and confidence-level sensitivity.
\item Implement the fault-injection state machine and verify identical agent-visible prefixes for the paired ambiguity histories. Account for pending requests and acknowledgments through receipt, processing, acknowledgment, crash, and checkpoint boundaries.
\item Validate actual external-effect atomicity, durability, payload binding, stable identity, concurrency, failover, retention, and downstream coverage. A local key table or transactional dispatch intention is not sufficient evidence.
\item Test queued copies, paused-worker resumption, record expiry, garbage-collection ordering, immutable authorization, and universal execution-time gating before claiming safety under finite retention.
\item Validate epoch advancement and restart behavior, including lost or delayed replies and suppression of abandoned controller flows. Establish progress premises separately from at-most-once safety.
\item Implement an independently inspectable outcome observer with complete effect and report capture, cutoff ordering, namespace isolation, conflict handling, and explicit nonexecution evidence. Keep mapper conformance distinct from reference validity.
\item Run matched-budget, task-blocked comparisons of the specified recovery policies with randomized condition order and independently audited outcomes. Report uncertainty, severity, safe stopping, completion, recovery, latency, and monetary or compute cost.
\item Extend the evaluation to held-out tasks, changing environments, dependent workflows, delayed harms, and non-ledger outcomes before broadening the reliability claim.
\end{itemize}
\section{Frozen obligation reconciliation}
Unlisted obligations are not presumed resolved. The archived research snapshot retains complete source references and the full reconciliation.
\begin{itemize}
\item Implement and audit the fault-injection harness, especially the assertion that the two histories expose identical information to the agent while the evaluator can distinguish their external states. \textit{ (Not listed by author; not presumed resolved; 10 saved source reference(s))}
\item Measure whether real agent policies preserve operation identity and payload across retries, checkpoint restoration, replanning, and concurrent workers. \textit{ (Not listed by author; not presumed resolved; 10 saved source reference(s))}
\item Test atomicity, crash durability, key retention, collision prevention, and payload-mismatch handling in any service claimed to satisfy the sufficient-repair assumptions. \textit{ (Not listed by author; not presumed resolved; 10 saved source reference(s))}
\item Estimate deployment frequencies and asymmetric costs of missing versus duplicate effects before choosing a practical retry policy; the worst-case bound does not supply these quantities. \textit{ (Not listed by author; not presumed resolved; 6 saved source reference(s))}
\item Evaluate authoritative lookup, escalation, compensation, and safe termination as additional policies, distinguishing harm containment from successful and verified task completion. \textit{ (Not listed by author; not presumed resolved; 10 saved source reference(s))}
\item Extend beyond one external action to dependent workflows with stale reads and partial multi-service commits; the current proof does not establish end-to-end workflow reliability. \textit{ (Not listed by author; not presumed resolved; 6 saved source reference(s))}
\item Run repeated, task-blocked comparisons with independent outcome audits and uncertainty estimates. Scripted protocol agreement alone would not establish general agent reliability. \textit{ (Not listed by author; not presumed resolved; 3 saved source reference(s))}
\item For open-ended tasks without an authoritative effect ledger, develop and validate an outcome oracle rather than transferring this experiment's observability assumptions. \textit{ (Not listed by author; not presumed resolved; 8 saved source reference(s))}
\item Implement and audit the fault-injection harness, including equality of all agent-accessible information and separation of evaluator-only external-state evidence. \textit{ (Not listed by author; not presumed resolved; 4 saved source reference(s))}
\item Measure operation-key and payload preservation by real agent policies across retry, checkpoint restoration, replanning, and concurrent execution. \textit{ (Not listed by author; not presumed resolved; 4 saved source reference(s))}
\item Test service atomicity, crash durability, retention, collision prevention, and payload-mismatch behavior. \textit{ (Not listed by author; not presumed resolved; 4 saved source reference(s))}
\item Estimate ambiguous-commit frequencies and asymmetric missing-effect and duplicate-effect costs before recommending a deployment policy. \textit{ (Not listed by author; not presumed resolved; 9 saved source reference(s))}
\item Evaluate authoritative lookup, escalation, compensation, and safe termination while separating harm containment from successful, verified completion. \textit{ (Not listed by author; not presumed resolved; 4 saved source reference(s))}
\item Extend the analysis to dependent workflows, stale observations, and partial multi-service commits. \textit{ (Not listed by author; not presumed resolved; 4 saved source reference(s))}
\item Conduct repeated, task-blocked comparisons with independent outcome audits and appropriate uncertainty estimates. \textit{ (Not listed by author; not presumed resolved; 9 saved source reference(s))}
\item Develop and validate outcome oracles for tasks without an authoritative external-effect ledger. \textit{ (Not listed by author; not presumed resolved; 4 saved source reference(s))}
\item Material, recoverable evidence gap: identical agent observations and independently observable external effects are assumptions of the proof, not properties demonstrated for an implemented harness. Timing, tool metadata, restored memory, and outstanding requests could invalidate the intended experimental pairing. \textit{ (Not listed by author; not presumed resolved; 7 saved source reference(s))}
\item Material, recoverable evidence gap: the sufficient repair assumes atomic, durable coupling of the business effect and deduplication record. A service that merely maintains a key table has not thereby satisfied that assumption. \textit{ (Not listed by author; not presumed resolved; 2 saved source reference(s))}
\end{itemize}
682 further distinct obligations are retained in the archived reconciliation.
\section{Conclusions}
A defensible reliability claim must identify its execution population, resource limits, observation horizon, and actual success rubric. Repeated runs and paired tasks help define an evaluation, but their dependence structure and weighting remain part of the claim.

Ambiguous external writes reveal why local completion evidence is insufficient. The corrected analytical lesson is precise: retry issuance is not commitment, and different effect-marker orderings threaten different properties.

Recovery safeguards should be evaluated as contracts over actual effects, identity, persistence, and execution paths. At-most-once safety, timely delivery, truthful reporting, and containment require separate judgments.

The strongest supported outcome is a conditional evaluation framework and a set of discriminating tests, not evidence that a particular agent architecture improves reliability. Demonstrating that improvement requires independently validated outcome evidence and controlled repeated comparisons.

\textit{Claims discussed: }C31, C32, C35, C36, C38, C40, C41, C42, C44, C45, C47
\appendix
\section{Selected source-linked research excerpts}
The following are bounded, selected excerpts of original research records. They are not exhaustive; consult the full frozen snapshot archived with this paper. They do not constitute independent verification.
\subsection{Research excerpt 1}
\textit{Source and field: }[Research session 6, artifact 4]: skeptic\_output/reasoning\_assessments\par
[\{"basis\_artifact\_ids": ["[Research session 6, artifact 2]", "[Research session 5, artifact 2]", "[Research session 6, artifact 3]", "[Research session 6, artifact 1]"], "disposition": "accepted", "objection": "", "rationale": "The supplied cycle 1 argument explicitly requires records to cover every possible retry and supplies the non-atomic and regenerated-key counterexamples. The peer architecture extract explicitly excludes expiry, while the measurement peer concerns a frozen roster.

None establishes finite-retention safety, and the current proof does not transfer their unsupported claims.", "step\_number": 1\}, \{"basis\_artifact\_ids": [], "disposition": "accepted", "objection": "", "rationale": "These are explicit premises of the conditional claim. Atomic checking with the actual effect, per-key ordering of deletion, crash durability, and monotonic retention time make the later concurrency and expiry arguments well-defined. They are not evidence that any deployed interface satisfies this model.", "step\_number": 2\}, \{"basis\_artifact\_ids": [], "disposition": "accepted", "objection": "", "rationale": "The first commit at 1 gives expiry at 11.

Deletion at 11 and an outstanding request's transition at 12 permit a second commit under the specified absent-record rule. Stopping sends at 3 neither removes nor disables that request. Thus C=2 follows without violating any expiry-only premise.", "step\_number": 3\}, \{"basis\_artifact\_ids": [], "disposition": "accepted", "objection": "", "rationale": "A request admitted at 2.5 can remain queued until 12, so an ingress\ldots{}

\subsection{Research excerpt 2}
\textit{Source and field: }[Research session 6, artifact 5]: skeptic\_output/reasoning\_assessments\par
[\{"basis\_artifact\_ids": ["[Research session 6, artifact 3]", "[Research session 5, artifact 2]", "[Research session 6, artifact 2]", "[Research session 6, artifact 1]"], "disposition": "accepted", "objection": "", "rationale": "The supplied cycle-1 argument explicitly requires nonexpiring receipts, so removing that premise requires a new invariant. The peer extracts likewise provide no finite-retention proof. The current ledger is nonempty, and the earlier wording error is explicitly corrected without claiming historical resolution.", "step\_number": 1\}, \{"basis\_artifact\_ids": [], "disposition": "accepted", "objection": "", "rationale": "These are explicit premises of the conditional claim.

Logical expiration at t=c+T, no replay refresh, and atomic actual-effect/receipt commitment specify the boundary behavior needed for the trace and proofs; none is presented as a demonstrated service property.", "step\_number": 2\}, \{"basis\_artifact\_ids": [], "disposition": "accepted", "objection": "", "rationale": "The delayed-copy and restart premises explicitly permit processing long after sending while preserving committed state.

Requiring eligibility and deduplication checks at the effect-producing transition excludes stale authorization from the conditional model.", "step\_number": 3\}, \{"basis\_artifact\_ids": [], "disposition": "accepted", "objection": "", "rationale": "At t=12 the receipt committed at t=1 is invalid because its expiry is 11. The ordinary service therefore takes its absent-receipt execution branch and increments C from one to two. Stable keys, atomic\ldots{}

\subsection{Research excerpt 3}
\textit{Source and field: }[Research session 7, artifact 1]: skeptic\_output/reasoning\_assessments\par
[\{"basis\_artifact\_ids": ["[Research session 6, artifact 1]"], "disposition": "accepted", "objection": "", "rationale": "The supplied extract of [Research session 6, artifact 1] is explicitly incomplete. Using it to identify recorded limitations while independently deriving the current result avoids importing unavailable proof steps. No retry-commitment inference is used.", "step\_number": 1\}, \{"basis\_artifact\_ids": [], "disposition": "accepted", "objection": "", "rationale": "Accepted as a proposed protocol premise, not an observed evaluation. Four schedules repeated four times define 16 pairs, and retaining all 32 scheduled executions with prespecified nonexecution failure keeps the denominator fixed.

The explicitly descriptive contrast does not assert a causal effect.", "step\_number": 2\}, \{"basis\_artifact\_ids": [], "disposition": "accepted", "objection": "", "rationale": "Accepted under the explicitly listed sound-certificate and well-defined-cutoff premises. The conjunction Y=E AND R distinguishes actual external success from reporting success. The step correctly treats certificate soundness as assumed rather than established by separate retention.", "step\_number": 3\}, \{"basis\_artifact\_ids": [], "disposition": "accepted", "objection": "", "rationale": "If each fact set contains its actual binary fact, the image under AND contains actual Y. Taking a Cartesian product preserves paired containment without component independence.

Widening conflicting facts avoids unsound empty intersections; a trusted duplicate state and a success report correctly imply failure\ldots{}

\subsection{Research excerpt 4}
\textit{Source and field: }[Research session 4, artifact 1]: skeptic\_output/reasoning\_assessments\par
[\{"basis\_artifact\_ids": [], "disposition": "accepted", "objection": "", "rationale": "Uniform selection within a fixed r-run batch defines a coherent marginal execution population. The localized unconditional error probability c differs correctly from a conditional reference-error rate. The allowance is accepted only as an explicit premise of the conditional claim.", "step\_number": 1\}, \{"basis\_artifact\_ids": [], "disposition": "accepted", "objection": "", "rationale": "Partitioning accepted failures by R gives Pr(O=1,Y=0)$\leq$a+c$\leq$a+η.

Accepted true successes form a subset of all true successes, so subtraction from q and clipping at zero establish the bound without judge independence.", "step\_number": 2\}, \{"basis\_artifact\_ids": [], "disposition": "accepted", "objection": "", "rationale": "For any feasible observed joint distribution, setting truth to failure outside the jointly accepted cell and allocating min(s,η) failure mass inside it attains max(0,s-η). Universal success is also feasible.

Mixtures preserve observations and the inequality on c, establishing every intermediate reliability value.", "step\_number": 3\}, \{"basis\_artifact\_ids": [], "disposition": "accepted", "objection": "", "rationale": "The two worlds have identical observed judge outputs but different truth labels and reliabilities. This establishes nonidentification from agreement alone, not an empirical benchmark failure rate or a contradiction to a separately justified η\textless{}1.", "step\_number": 4\}, \{"basis\_artifact\_ids": [], "disposition": "accepted", "objection": "", "rationale": "Averaging the run-specific\ldots{}

\subsection{Research excerpt 5}
\textit{Source and field: }[Research session 5, artifact 3]: skeptic\_output/reasoning\_assessments\par
[\{"basis\_artifact\_ids": [], "disposition": "accepted", "objection": "", "rationale": "As a premise of the explicitly conditional claim, fixed strata and weights together with uniform selection among scheduled runs define a coherent task-weighted estimand. Isolated paired copies and unrestricted within-cluster dependence do not conflict with independence between complete task clusters.", "step\_number": 1\}, \{"basis\_artifact\_ids": ["[Research session 4, artifact 1]"], "disposition": "accepted", "objection": "", "rationale": "The listed assumptions explicitly include these unconditional directional error bounds.

Intervention false certification can inflate intervention reliability, while baseline missed successes can depress baseline certification; these are the relevant adverse directions for a lower improvement bound. They are premises, not established calibration results. The supplied prior extract discusses only the lower-bound-oriented allowance.", "step\_number": 2\}, \{"basis\_artifact\_ids": [], "disposition": "accepted", "objection": "", "rationale": "The binary identity is exact. Dropping intervention missed successes and baseline false certifications gives the stated lower and upper bounds respectively.

Subtraction and nonnegative weighted summation preserve the inequality without judge independence.", "step\_number": 3\}, \{"basis\_artifact\_ids": [], "disposition": "accepted", "objection": "", "rationale": "With S\_A=1, S\_B=0, and both truth labels equal to one, the certification lower bounds differ by one while actual improvement is zero. Baseline missed-success\ldots{}

\subsection{Research excerpt 6}
\textit{Source and field: }[Research session 6]: steps/8/output/reasoning\_assessments\par
[\{"basis\_artifact\_ids": ["[Research session 6, artifact 2]", "[Research session 5, artifact 2]", "[Research session 6, artifact 3]", "[Research session 6, artifact 1]"], "disposition": "accepted", "objection": "", "rationale": "The supplied cycle 1 argument explicitly requires records to cover every possible retry and supplies the non-atomic and regenerated-key counterexamples. The peer architecture extract explicitly excludes expiry, while the measurement peer concerns a frozen roster.

None establishes finite-retention safety, and the current proof does not transfer their unsupported claims.", "step\_number": 1\}, \{"basis\_artifact\_ids": [], "disposition": "accepted", "objection": "", "rationale": "These are explicit premises of the conditional claim. Atomic checking with the actual effect, per-key ordering of deletion, crash durability, and monotonic retention time make the later concurrency and expiry arguments well-defined. They are not evidence that any deployed interface satisfies this model.", "step\_number": 2\}, \{"basis\_artifact\_ids": [], "disposition": "accepted", "objection": "", "rationale": "The first commit at 1 gives expiry at 11.

Deletion at 11 and an outstanding request's transition at 12 permit a second commit under the specified absent-record rule. Stopping sends at 3 neither removes nor disables that request. Thus C=2 follows without violating any expiry-only premise.", "step\_number": 3\}, \{"basis\_artifact\_ids": [], "disposition": "accepted", "objection": "", "rationale": "A request admitted at 2.5 can remain queued until 12, so an ingress\ldots{}

\subsection{Research excerpt 7}
\textit{Source and field: }[Research session 6, artifact 3]: skeptic\_output/reasoning\_assessments\par
[\{"basis\_artifact\_ids": ["[Research session 5, artifact 2]"], "disposition": "challenged", "objection": "Nonmaterial factual correction: five current ledger entries are supplied, so 'empty current ledger' is inaccurate. No historical resolution should be inferred either way; no rewrite cycle is needed for this wording alone.", "rationale": "The source is explicitly an incomplete extract, and declining to import its unshown proof is appropriate. Omitted history cannot establish resolution.

The only challenged component is the description of the current ledger.", "step\_number": 1\}, \{"basis\_artifact\_ids": ["[Research session 5, artifact 2]"], "disposition": "accepted", "objection": "", "rationale": "Under identical pre-retry information and assured retry commitment, success in the zero-effect history requires retry, while success in the one-effect history requires no retry. The stated bounds follow.

With every retry dropped, q=1 but the second history retains exactly one effect, directly disproving the relaxed second inequality.", "step\_number": 2\}, \{"basis\_artifact\_ids": [], "disposition": "accepted", "objection": "", "rationale": "These are explicit premises of the conditional claim, not established properties of deployed tools. They specify the durable states, replay behavior, serialization, operation identity, and excluded repair mechanisms sufficiently to determine the following traces.", "step\_number": 3\}, \{"basis\_artifact\_ids": [], "disposition": "accepted", "objection": "", "rationale": "Each table cell follows by applying the defined replay rule to the crash\ldots{}

\subsection{Research excerpt 8}
\textit{Source and field: }[Research session 7]: outcome/solution\_review/strategy\_findings/1/reasoning\_steps\par
[\{"argument": "Let C(k,t) count historical committed primary effects attributed to logical operation k; it is monotone and is not reduced by refunds or corrective records. The present incident domain permits a nonzero or corrupted initial state, unlike the clean-state progress domain described in the supplied [Research session 7, artifact 4] extract. A trustworthy post-fence read C\textgreater{}1 therefore exits that earlier guarantee's premises; it does not refute a result conditional on clean initialization. The original [Research session 7, artifact 2] is pointer-only here and has not been inspected.

Its reported invariant is not a premise of the new containment proof.", "basis\_artifact\_ids": ["[Research session 7, artifact 4]"], "epistemic\_status": "deduction", "number": 1, "research\_artifact\_id": "[Research session 7, artifact 7]", "review": \{"basis\_artifact\_ids": [], "disposition": "accepted", "objection": "", "rationale": "Defining C as historical committed effects makes compensation non-subtractive. A trustworthy duplicate count indicates that at least one premise of a clean-state safety guarantee is inapplicable; it need not specifically identify initialization as the failed premise.

The supplied ledger supports this scope comparison, not verification of the unavailable earlier proof."\}, "title": "Scope and provenance"\}, \{"argument": "Introduce Close(k,incident\_id), a proposed administrative primitive, not an existing property inferred from Advance. Close durably changes the operation's gate to CLOSED at a linearization point q. Every covered primary effect\ldots{}

\subsection{Research excerpt 9}
\textit{Source and field: }[Research session 7, artifact 5]: skeptic\_output/reasoning\_assessments\par
[\{"basis\_artifact\_ids": ["[Research session 7, artifact 3]", "[Research session 6, artifact 1]"], "disposition": "accepted", "objection": "", "rationale": "The supplied cycle-1 extract contains both the valid invariant and the challenged split-commit assertion. The current step explicitly withdraws the overstatement and does not import the truncated historical audit theorem.", "step\_number": 1\}, \{"basis\_artifact\_ids": [], "disposition": "accepted", "objection": "", "rationale": "These are proposed premises, not inspected implementation facts.

The E schema permits duplicate request identifiers, D supplies logical-key uniqueness, and the text explicitly identifies linked-row consistency as an additional application obligation rather than falsely attributing it to the foreign key.", "step\_number": 2\}, \{"basis\_artifact\_ids": [], "disposition": "accepted", "objection": "", "rationale": "Under the conditional claim, atomic rollback includes E, ambiguous commits are resolved through a new D lookup, and restart never blindly resumes an old append.

These requirements close the specified duplicate-creation paths without asserting an actual database satisfies them.", "step\_number": 3\}, \{"basis\_artifact\_ids": ["[Research session 7, artifact 3]"], "disposition": "accepted", "objection": "", "rationale": "The first successful absent-key transaction creates both rows. D's unique key forces competing transactions to abort their E writes, and retained D prevents later insertion. Complete-state crash recovery preserves the invariant; a safety violation in an infinite execution\ldots{}

\subsection{Research excerpt 10}
\textit{Source and field: }[Research session 7, artifact 7]: skeptic\_output/reasoning\_assessments\par
[\{"basis\_artifact\_ids": [], "disposition": "accepted", "objection": "", "rationale": "Defining C as historical committed effects makes compensation non-subtractive. A trustworthy duplicate count indicates that at least one premise of a clean-state safety guarantee is inapplicable; it need not specifically identify initialization as the failed premise. The supplied ledger supports this scope comparison, not verification of the unavailable earlier proof.", "step\_number": 1\}, \{"basis\_artifact\_ids": [], "disposition": "accepted", "objection": "", "rationale": "This is an explicit premise of the conditional claim.

Atomic ordering of closure with every covered effect creation, persistent CLOSED state, and prohibition of ordinary reopening are sufficient contract ingredients for the later invariant; no implementation is inferred.", "step\_number": 2\}, \{"basis\_artifact\_ids": [], "disposition": "accepted", "objection": "", "rationale": "These are additional stated authority, persistence, and evidence premises, not established capabilities. Durable fail-closed dispatch suppression supports restart behavior, while certificate identity and ordering requirements prevent confusing a requested close with a confirmed close.

Outbox persistence alone is correctly not treated as delivery.", "step\_number": 3\}, \{"basis\_artifact\_ids": [], "disposition": "accepted", "objection": "", "rationale": "The proposed incident action suppresses new autonomous business dispatches without claiming cancellation of already pending effects. Its distinction between immediate local stopping and suppression\ldots{}

\subsection{Research excerpt 11}
\textit{Source and field: }[Research session 7, artifact 9]: skeptic\_output/reasoning\_assessments\par
[\{"basis\_artifact\_ids": [], "disposition": "challenged", "objection": "[Research session 6, artifact 1] is supplied only as a pointer, not the substantive extract described in this step. Use the ledger as a recorded constraint, not as an inspected proof.", "rationale": "The continuity restriction and preservation of unknown duties are appropriate, but the claimed inspection exceeds the supplied evidence.", "step\_number": 1\}, \{"basis\_artifact\_ids": [], "disposition": "accepted", "objection": "", "rationale": "This is an explicit standalone rubric, conditional on the listed well-defined isolated history premise.

Exact-one commit and exact-one valid final receipt make additional events failures, and historical before-H labels are not rewritten by later events.", "step\_number": 2\}, \{"basis\_artifact\_ids": [], "disposition": "accepted", "objection": "", "rationale": "The specified branches preserve nonempty component sets. Completeness is required for exact-one positive conclusions, while correctly scoped pre-H irreversible violations can establish failure without complete history.

Conflict widening is conservative rather than an empty-set inference.", "step\_number": 3\}, \{"basis\_artifact\_ids": [], "disposition": "accepted", "objection": "", "rationale": "M compares output with the input-version specification; V compares certificate assertions with history; Z tests membership of the actual outcome; B depends on review-time trust information. These definitions do not equate ignorance with truth or a hold with task failure.", "step\_number": 4\}, \{"basis\_artifact\_ids": [],\ldots{}

\subsection{Research excerpt 12}
\textit{Source and field: }[Research session 1]: outcome/solution\_review/proof\_obligations\par
["Analytically discharged within the conditional model, cycle 1, current Investigator reasoning entries 1\textendash{}3: construct indistinguishable histories and establish the sharp minimax error bound.", "Analytically discharged within the conditional model, cycle 1, current Investigator reasoning entry 4: distinguish conditional average error from worst-case error and overall deployment failure.", "Analytically discharged within the conditional model, cycle 1, current Investigator reasoning entries 5\textendash{}6: explain why local rollback is insufficient and why atomic durable same-key deduplication preserves exactly one effect.", "Open: establish that a concrete harness implements the assumed information boundary and transition semantics.", "Open: establish that concrete agent policies and services satisfy the repair assumptions.", "Open: measure controlled agent-level outcomes and establish the applicability limits for deployment and multi-step workflows."]

\subsection{Research excerpt 13}
\textit{Source and field: }[Research session 1, artifact 1]: skeptic\_output/reasoning\_assessments\par
[\{"basis\_artifact\_ids": [], "disposition": "accepted", "objection": "", "rationale": "The uniform run selection makes E[Y\_J] the expected within-batch success fraction under the fixed protocol. Explicitly including crashes and timeouts avoids conditioning on successful execution or oracle acceptance. This defines an estimand rather than establishing deployment representativeness.", "step\_number": 1\}, \{"basis\_artifact\_ids": [], "disposition": "accepted", "objection": "", "rationale": "Partitioning the joint binary outcomes gives q=P(Y=1,Z=1)+f and p=P(Y=1,Z=1)+g, so p=q-f+g.

The two constant-oracle examples correctly demonstrate nonidentification from q alone, and f is correctly distinguished from a conditional false-positive rate.", "step\_number": 2\}, \{"basis\_artifact\_ids": [], "disposition": "accepted", "objection": "", "rationale": "Under the listed independent-cluster assumptions, each X\_i lies in [0,1] and has expectation q because uniform run selection averages the r oracle decisions. One-sided Hoeffding therefore has exponent 2Nt\textasciicircum{}2.

Dependence among runs within X\_i does not invalidate that application or justify an Nr exponent.", "step\_number": 3\}, \{"basis\_artifact\_ids": [], "disposition": "accepted", "objection": "", "rationale": "Distribution-matched independent audit clusters, uniform run selection, and correct adjudication make the A\_l independent Bernoulli(f). For k\textless{}m, the binomial lower-tail probability decreases continuously from 1 to 0 as u increases, giving the stated unique root. Inversion yields the conservative one-sided upper confidence limit, including\ldots{}

\subsection{Research excerpt 14}
\textit{Source and field: }[Research session 1, artifact 2]: objective\_review/proof\_obligations\par
["Analytically discharged within the conditional model, cycle 1, current Investigator reasoning entries 1\textendash{}3: construct indistinguishable histories and establish the sharp minimax error bound.", "Analytically discharged within the conditional model, cycle 1, current Investigator reasoning entry 4: distinguish conditional average error from worst-case error and overall deployment failure.", "Analytically discharged within the conditional model, cycle 1, current Investigator reasoning entries 5\textendash{}6: explain why local rollback is insufficient and why atomic durable same-key deduplication preserves exactly one effect.", "Open: establish that a concrete harness implements the assumed information boundary and transition semantics.", "Open: establish that concrete agent policies and services satisfy the repair assumptions.", "Open: measure controlled agent-level outcomes and establish the applicability limits for deployment and multi-step workflows."]

\subsection{Research excerpt 15}
\textit{Source and field: }[Research session 4]: outcome/solution\_review/proof\_obligations\par
["Lifecycle and request conservation: cycle 2, steps 1\textendash{}4 establish explicit location transfers and crash outcomes under the chosen transport premises.", "Durable atomic safety: cycle 2, step 5 establishes preservation under the refined transitions; cycle 1, step 3 in artifact [Research session 4, artifact 2] supplies the compatible saved inductive argument.", "Affected finite traces: cycle 2, steps 6\textendash{}10 supply consistent per-event T5, T8, and both T12 tables, including consumed deliveries and acknowledgments.", "Safety versus existence: cycle 1, step 4 in artifact [Research session 4, artifact 2] requires a completed matching PROCESS for exactly-one existence; cycle 2 adds no liveness guarantee.", "Naive split-write failure boundary: cycle 1, step 7 in artifact [Research session 4, artifact 2] retains valid local duplicate-effect and false-completion counterexamples; these do not refute all non-atomic recovery protocols.", "Model-to-system correspondence, real-agent identity preservation, controlled causal benefit, broader workflow applicability, and historical reconciliation remain unestablished."]

\subsection{Research excerpt 16}
\textit{Source and field: }[Research session 2]: outcome/conclusion\par
"\#\#\# Unresolved investigation\textbackslash{}n\textbackslash{}nThe session is queued and has not saved a complete Investigator/Skeptic finding. No answer to the project question is established yet.\textbackslash{}n\textbackslash{}n\#\#\# Findings across strategies\textbackslash{}n- **Develop a rigorous framework for measuring AI agent reliability \textemdash{} unresolved.**\textbackslash{}n- **Understand failure propagation and recovery in autonomous AI agents \textemdash{} unresolved.**\textbackslash{}n- **Test architectural mechanisms for dependable agent outcomes \textemdash{} unresolved.**\textbackslash{}n\textbackslash{}n\#\#\# Session caveats\textbackslash{}n- No complete research artifact is available.\textbackslash{}n\textbackslash{}n\#\#\# Verification and scope\textbackslash{}n\textbackslash{}nVerification: Incomplete.\textbackslash{}n\textbackslash{}nThis Conclusion summarizes agent-reported findings for this session.

SIGMA has not independently verified the result; inspect the detailed Outcome and saved evidence before relying on it."

4 further research records are retained in the archived source snapshot.
\section*{Evidence and authorship note}
Snapshot SHA-256: 080d3b70b67cfbe346acdb0f565d4a61c09ec4819229f2dd92dded43d9de40d9
Generated by an AI author from the cited saved research snapshot. This document is not independent verification. No external bibliography was added.
\section*{Validation warnings}
\begin{itemize}
\item Evidence indexing omitted 32185 material saved research field fragment(s). Those fragments were not reviewed for this paper; the interpretation is provisional where the omitted evidence may matter. Consult the frozen source snapshot and omission ledger for full details.
\item Author self-review: the bounded preview does not expose complete statistical proofs or enough numerical detail to reproduce a confidence bound. No omitted formula has been inferred.
\item The proposed retry bound was explicitly corrected: without assured retry commitment, an always-retry policy with lost retries defeats the claimed upper bound in the already-committed history.
\item The blanket claim that split commits invalidate at-most-once safety was explicitly withdrawn. Effect-first duplication and marker-first loss of delivery must remain distinct.
\item A mapper that follows its specification is not automatically defective when its input evidence excludes the actual outcome. Reference validity and mapper conformance require separate labels.
\item Service invariants, sound certificates, and possible-outcome containment are conditional premises, not observed implementation properties. The treatment invariant must not be reused as the outcome oracle.
\item Preflight and context-assembly failures in the research records do not provide controlled evidence of deployed-agent failure prevalence or architectural effectiveness.
\item The selection contains no genuine external bibliographic sources sufficient for a literature review or bibliography. Benchmark sensitivity and contamination remain research questions rather than findings here.
\item The planned qualitative comparison table must be populated only from inspected mechanisms and explicit assumptions; it must not imply measured reliability, latency, or cost rankings.
\item The evidence index explicitly omitted 32189 saved field fragment(s); see the generation coverage ledger and review the complete archived snapshot.
\item Structured paper indexing selected 441 source-linked text fragments from 36 archived sources without AI extraction calls. The complete frozen snapshot remains archived; this bounded index does not constitute review of every saved field.
\item C31 is labeled established, but cited saved evidence includes unverified, incomplete, or conjectural material.
\item C31 is labeled established while the frozen evidence contains counterexample or contradiction language; inspect the cited sources before accepting it.
\item C32 is labeled established, but cited saved evidence includes unverified, incomplete, or conjectural material.
\item C32 is labeled established while the frozen evidence contains counterexample or contradiction language; inspect the cited sources before accepting it.
\item C35 is labeled established, but cited saved evidence includes unverified, incomplete, or conjectural material.
\item C35 is labeled established while the frozen evidence contains counterexample or contradiction language; inspect the cited sources before accepting it.
\item C36 is labeled established, but cited saved evidence includes unverified, incomplete, or conjectural material.
\item C36 is labeled established while the frozen evidence contains counterexample or contradiction language; inspect the cited sources before accepting it.
\item C37 is labeled established, but cited saved evidence includes unverified, incomplete, or conjectural material.
\item A diagram's graph could not be laid out without likely crossings or overlap; its complete nodes, boundaries, and relationships are shown in a readable text representation.
\end{itemize}
\section*{Typesetting verification note}
Needs review: TeX verification is incomplete. The authored research content and safe accepted typesetting corrections were retained. Review the saved proofreading diagnostics before relying on the layout.
\end{document}
